% Part: first-order-logic
% Chapter: natural-deduction
% Section: rules-and-proofs

\documentclass[../../../include/open-logic-section]{subfiles}

\begin{document}

\iftag{FOL}
      {\olfileid{fol}{ntd}{rul}}
      {\olfileid{pl}{ntd}{rul}}

\olsection{قاعدې او \usetoken{P}{derivation}}

\begin{explain}
د فطري استنتاج له نظامونو څخه موخه دا ده چې په رياضيکي ثبوت کښې له
کارېدونکي غيررسمي استدلال سره نژدې موازي وي (همدا وجه ده چې تر يوه
حده “فطري” دے)۔ د فطري استنتاج ثبوتونه له فرضيو پيلېږي۔ بيا د
استنتاج قاعدې ورباندې کارول کېږي۔ فرضيې د~\Intro{\lnot}،
\Intro{\lif}، \iftag{FOL}{\Elim{\lor} او \Elim{\lexists}}{ او
  \Elim{\lor}} د استنتاج د قاعدو په وسيله “!!{discharged}” کېږي،
او د روښانتيا دپاره د !!{discharged} شوې فرضيې نښان د استنتاج تر
څنګ ليکل کېږي۔
\end{explain}

\begin{defn}[فرضيه]
يوه \emph{فرضيه} هره هغه !!{sentence} ده چې د هرې څانګې په تر ټولو
پورتني ځاے کښې وي۔
\end{defn}

په فطري استنتاج کښې !!^{derivation}s د !!{sentence}s ځانګړې ونې دي؛
تر ټولو پورتنۍ !!{sentence}s فرضيې وي، او که يوه !!a{sentence} تر
يوې، دوو يا دريو نورو جملو لاندې وي، بايد په سمه توګه د استنتاج له
يوې قاعدې څخه پايله شي۔ د استنتاج په سر کښې !!{sentence}s د هغې
\emph{مقدمې} او لاندې !!{sentence} د استنتاج \emph{پايله} بلل کېږي۔
قاعدې په جوړو کښې راځي: د هر !!{operator} دپاره يوه د داخلولو او
يوه د ايستلو قاعده۔ دوئ په پايله کښې !!a{operator} داخلوي يا
!!a{operator} د قاعدې له يوې مقدمې څخه باسي۔ ځينې قاعدې اجازه
ورکوي چې د يوه ټاکلي
ډول فرضيه \emph{!!{discharged}} شي۔ د دې د ښودلو دپاره چې کومه
فرضيه د کوم استنتاج په وسيله !!{discharged} کېږي، فرضيې او استنتاج
دواړو ته نښانونه ورکوو۔ دا د فرضيې په “$\Discharge{!A}{n}$” ليکلو
ښودل کېږي۔ \textbf{د ژباړې د نسخې سمون (OLND-001):} سرچينه په دې
ونه کښې لاندې ولاړ توکي “سېکوېنټونه” بولي؛ دلته “جملې” راوړل شوې،
ځکه خپله سرچينه ونه د جملو ونه تعريفوي او قاعدې پر جملو لګېږي۔

% Only include this sentence is on of \land, lor, lif or lnot is defined.

\iftag{notprvNot,notprvAnd,notprvOr,notprvIf}{}%
	{دا دود دے چې د ټولو !!{operator}s~$\land$، $\lor$، $\lif$،
    $\lnot$ او~$\lfalse$ دپاره قاعدې په پام کښې ونيول شي، که څه هم
    ځينې ئې تعريف شوي وي۔}



\end{document}
